\section*{Висновки до третього роздылу}

На основі проведеної програмної реалізації та детального аналізу отриманих результатів можна зробити наступні висновки:

\begin{itemize}
	\item \textbf{Створення середовища для тестування:} Розроблено гнучкий модуль автоматизованого тестування (\texttt{runBenchmark}), який дозволяє проводити експерименти як на окремих зображеннях, так і здійснювати пакетну обробку цілих директорій. Це дозволило зібрати статистику на великих масивах даних.
	
	\item \textbf{Об'єктивність оцінювання:} Для коректного порівняння методів було успішно застосовано методологію інверсного моделювання, що забезпечило абсолютно рівні початкові умови для тестування алгоритмів. Оцінка якості проводилася за загальноприйнятими математичними метриками: середньоквадратичною похибкою MSE, піковим відношенням сигналу до шуму (PSNR) та індексом структурної подібності (SSIM).
	
	\item \textbf{Загальна конкурентоспроможність:} Згідно з представленими результатами тестування для різних датасетів та для різних коефіцієнтів масштабування ($s=2, 3, 4$) в обох напрямках обробки, прослідковується не погана конкурентоспроможність розробленого барицентричного методу.
	
	\item \textbf{Підтвердження апаратної оптимізації:} Порівняння часу виконання експериментів доводить критичну важливість розроблених інженерних рішень. Перехід до матричної форми з використанням бібліотеки OpenCV (\texttt{SBI Матричне множення з OpenCV}) дозволив кардинально прискорити обчислення порівняно з базовою реалізацією, зробивши метод придатним для практичного застосування.
\end{itemize}